\chapter{Wasserstein 距離 \texorpdfstring{$W_p$}{Wp}}
\label{ch:wasserstein-metrics}

\section{設定と定義}
\label{sec:wp-definition}

$(\mathcal{X},d)$ を Polish 空間（定義~\ref{def:meas-polish}）とする．
Borel $\sigma$-代数 $\Borel(\mathcal{X})$ 上の確率測度全体を
\[
  \Prob(\mathcal{X})
  :=
  \left\{
    \mu\colon\Borel(\mathcal{X})\to[0,1]
    \;\middle|\;
    \mu\text{ は測度},\;\mu(\mathcal{X})=1
  \right\}
\]
とおく．

\begin{definition}{有限 $p$ 次モーメントをもつ確率測度（Villani Def 6.4）}{wp-p-space}
  任意に点 $x_0\in \mathcal{X}$ を一つ固定する．$1\leq p<\infty$ に対し
  \[
    \Pp{p}(\mathcal{X})
    :=
    \left\{
      \mu\in\Prob(\mathcal{X})
      \;\middle|\;
      \int_\mathcal{X} d(x,x_0)^p\,\d\mu(x)<\infty
    \right\}
  \]
  と定める．積分 $\int_{\mathcal{X}}d(x,x_0)^p\,\d\mu(x)$ を
  $\mu$ の（$x_0$ を基点とする）\textbf{$p$ 次モーメント}という．
\end{definition}

\begin{lemma}{$\Pp{p}(\mathcal{X})$ の $x_0$ 非依存性（Villani Def 6.4）}{wp-p-indep}
  $1\leq p<\infty$ のとき，$\Pp{p}(\mathcal{X})$ は点 $x_0$ の選び方によらない．
  すなわち，任意の $x_0,x_0'\in \mathcal{X}$ と $\mu\in\Prob(\mathcal{X})$ に対して
  \[
    \int_\mathcal{X} d(x,x_0)^p\,\d\mu(x)<\infty
    \iff
    \int_\mathcal{X} d(x,x_0')^p\,\d\mu(x)<\infty
  \]
  が成り立つ．
\end{lemma}

\begin{proof}
  $x_0,x_0'\in \mathcal{X}$ と $\mu\in\Prob(\mathcal{X})$ を任意に取り，
  $\int_\mathcal{X} d(x,x_0)^p\,\d\mu(x)<\infty$ を仮定して，
  もう一方の積分が有限であることを示す．
  $a:=d(x,x_0)$，$b:=d(x_0,x_0')$ とおくと，三角不等式
  $d(x,x_0')\leq a+b$ と，$t\mapsto t^p$ が
  $[0,\infty)$ 上で凸であること（注意~\ref{rem:meas-real-facts}）から
  \[
    d(x,x_0')^p
    \leq (a+b)^p
    = 2^p\left(\frac{a}{2}+\frac{b}{2}\right)^p
    \leq 2^p\cdot\frac{a^p+b^p}{2}
    = 2^{p-1}\bigl(d(x,x_0)^p+d(x_0,x_0')^p\bigr).
  \]
  両辺を $\d\mu(x)$ で積分すると
  （右辺第2項は定数なので，$\mu(\mathcal{X})=1$ より積分値は $d(x_0,x_0')^p$）
  \[
    \int_\mathcal{X} d(x,x_0')^p\,\d\mu(x)
    \leq
    2^{p-1}\left(\int_\mathcal{X} d(x,x_0)^p\,\d\mu(x)+d(x_0,x_0')^p\right)
    <\infty
  \]
  を得る．$x_0$ と $x_0'$ を入れ替えれば逆向きの含意も同様に従う．
\end{proof}

\begin{definition}{Coupling（Villani Def 1.1）}{wp-coupling}
  距離空間 $\mathcal{X},\mathcal{Y}$ 上の確率測度
  $\mu\in\Prob(\mathcal{X})$，$\nu\in\Prob(\mathcal{Y})$ に対し
  \[
    \Couplings(\mu,\nu)
    :=
    \left\{
      \pi\in\Prob(\mathcal{X}\times\mathcal{Y})
      \;\middle|\;
      \begin{aligned}
        \pi(A\times\mathcal{Y})&=\mu(A) &&(\forall A\in\Borel(\mathcal{X})),\\
        \pi(\mathcal{X}\times B)&=\nu(B) &&(\forall B\in\Borel(\mathcal{Y}))
      \end{aligned}
    \right\}
  \]
  と定め，$\pi\in\Couplings(\mu,\nu)$ を $\mu$ と $\nu$ の
  \textbf{coupling} という．
\end{definition}

\begin{definition}{Wasserstein 距離（Villani Def 6.1）}{wp-wasserstein}
  $1\leq p<\infty$ と $\mu,\nu\in\Pp{p}(\mathcal{X})$ に対し
  \[
    \Wass_p(\mu,\nu)
    :=
    \inf_{\pi\in\Couplings(\mu,\nu)}
    \left(
      \int_{\mathcal{X}\times \mathcal{X}}d(x,y)^p\,\d\pi(x,y)
    \right)^{1/p}
  \]
  と定める．
\end{definition}

\begin{remark}{記法と可測性}{wp-welldef}
  距離関数 $d\colon \mathcal{X}\times \mathcal{X}\to[0,\infty)$ は，三角不等式から従う評価
  \[
    \abs{d(x,y)-d(x',y')}\leq d(x,x')+d(y,y')
  \]
  により連続，したがって Borel 可測である
  （注意~\ref{rem:meas-measurable-ops}）．
  $d\geq0$ より
  $\norm{d}_{L^p(\pi)}=\bigl(\int d^p\,\d\pi\bigr)^{1/p}$
  （$L^p$ ノルムは定義~\ref{def:meas-lp-norm}）だから，
  \[
    \Wass_p(\mu,\nu)
    =\inf_{\pi\in\Couplings(\mu,\nu)}\norm{d}_{L^p(\pi)}
  \]
  と書ける．
\end{remark}

\begin{lemma}{$\Wass_p$ の有限性（Villani Def 6.4）}{wp-finite}
  $1\leq p<\infty$，$\mu,\nu\in\Pp{p}(\mathcal{X})$ とすると，
  $\Wass_p(\mu,\nu)<\infty$ である．
\end{lemma}

\begin{proof}
  \textbf{1. coupling を一つ取る．}

  直積確率測度（定理~\ref{thm:meas-product}）
  $\pi_0:=\mu\otimes\nu$ を取る．任意の
  $A,B\in\Borel(\mathcal{X})$ に対し，
  \[
    \begin{aligned}
      \pi_0(A\times\mathcal{X})
      &=\mu(A)\nu(\mathcal{X})=\mu(A),\\
      \pi_0(\mathcal{X}\times B)
      &=\mu(\mathcal{X})\nu(B)=\nu(B).
    \end{aligned}
  \]
  したがって，coupling の定義（定義~\ref{def:wp-coupling}）より
  $\pi_0\in\Couplings(\mu,\nu)$ である．

  \textbf{2. $\pi_0$ に関する積分を $\mu,\nu$ に関する積分に直す．}

  $x_0\in\mathcal{X}$ を定義~\ref{def:wp-p-space}の基点とする．
  $d(x,x_0)^p$ は $x$ のみ，$d(x_0,y)^p$ は $y$ のみの
  非負可測関数である．$\pi_0=\mu\otimes\nu$ と Tonelli の定理
  （定理~\ref{thm:meas-product}）より
  \[
    \begin{aligned}
      &\int_{\mathcal{X}\times\mathcal{X}}
        d(x,x_0)^p\,\d\pi_0(x,y)\\
      &\quad=
      \int_{\mathcal{X}}
      \left(\int_{\mathcal{X}}d(x,x_0)^p\,\d\nu(y)\right)\d\mu(x)\\
      &\quad=
      \int_{\mathcal{X}}d(x,x_0)^p\,\nu(\mathcal{X})\,\d\mu(x)\\
      &\quad=
      \int_{\mathcal{X}}d(x,x_0)^p\,\d\mu(x)\\
      &\quad<\infty.
    \end{aligned}
  \]
  同様に，
  \[
    \begin{aligned}
      &\int_{\mathcal{X}\times\mathcal{X}}
        d(x_0,y)^p\,\d\pi_0(x,y)\\
      &\quad=
      \int_{\mathcal{X}}
      \left(\int_{\mathcal{X}}d(x_0,y)^p\,\d\mu(x)\right)\d\nu(y)\\
      &\quad=
      \int_{\mathcal{X}}\mu(\mathcal{X})\,d(x_0,y)^p\,\d\nu(y)\\
      &\quad=
      \int_{\mathcal{X}}d(x_0,y)^p\,\d\nu(y)\\
      &\quad<\infty.
    \end{aligned}
  \]
  ここで $\mu(\mathcal{X})=\nu(\mathcal{X})=1$ を用いた．
  また，各表示の最後の不等式は
  $\mu,\nu\in\Pp{p}(\mathcal{X})$ による．

  \textbf{3. 輸送コストを評価する．}

  任意の $(x,y)\in\mathcal{X}\times\mathcal{X}$ に対し，
  三角不等式より
  \[
    0\leq d(x,y)\leq d(x,x_0)+d(x_0,y).
  \]
  この不等式の両辺を $p$ 乗して $\pi_0$ で積分し，
  $1/p$ 乗を取る．さらに Minkowski の不等式
  （定理~\ref{thm:meas-minkowski}）を用いると，
  \[
    \begin{aligned}
      &\Bigl(\int_{\mathcal{X}\times\mathcal{X}}
        d(x,y)^p\,\d\pi_0(x,y)\Bigr)^{1/p}\\
      &\quad\leq
      \Bigl(\int_{\mathcal{X}\times\mathcal{X}}
        \bigl(d(x,x_0)+d(x_0,y)\bigr)^p\,\d\pi_0(x,y)\Bigr)^{1/p}\\
      &\quad\leq
      \Bigl(\int_{\mathcal{X}\times\mathcal{X}}d(x,x_0)^p\,\d\pi_0\Bigr)^{1/p}
      +
      \Bigl(\int_{\mathcal{X}\times\mathcal{X}}d(x_0,y)^p\,\d\pi_0\Bigr)^{1/p}\\
      &\quad=
      \Bigl(\int_{\mathcal{X}}d(x,x_0)^p\,\d\mu(x)\Bigr)^{1/p}
      +
      \Bigl(\int_{\mathcal{X}}d(x_0,y)^p\,\d\nu(y)\Bigr)^{1/p}\\
      &\quad<\infty.
    \end{aligned}
  \]
  2番目の不等式が Minkowski の不等式である．距離は非負だから，
  同定理に現れる絶対値は省略できる．また，等号と最後の不等式は
  2で確かめた計算による．

  \textbf{4. $\Wass_p$ の定義に戻る．}

  $\pi_0\in\Couplings(\mu,\nu)$ だから，$\Wass_p$ の定義と
  3の評価より
  \[
    0\leq\Wass_p(\mu,\nu)
    \leq
    \Bigl(\int_{\mathcal{X}\times\mathcal{X}}
      d(x,y)^p\,\d\pi_0(x,y)\Bigr)^{1/p}
    <\infty.
  \]
\end{proof}

\section{最適 coupling の存在}
\label{sec:wp-existence}

この節では，$\Wass_p$ の定義に現れる下限を実際に達成する coupling の存在を示す．
証明は，変分問題に対する直接法の次の四段階からなる．
\begin{enumerate}
  \item 輸送コストの下限に近づく最小化列 $(\pi_k)$ を取る．
  \item coupling の tightness と Prokhorov の定理を用いて，
    弱収束部分列 $\pi_{k_j}\Rightarrow\pi^*$ を取る．
  \item 弱極限 $\pi^*$ の二つの周辺分布が $\mu,\nu$ のままであることを示す．
  \item 輸送コストの下半連続性を用いて，$\pi^*$ が下限を達成することを示す．
\end{enumerate}
以下では，2に必要な tightness と，4に必要な下半連続性を先に準備する．
以下，積分域が明らかなときは
$\int_{\mathcal{X}\times\mathcal{X}}d(x,y)^p\,\d\pi(x,y)$ を
$\int d^p\,\d\pi$ と略記する．また，弱収束の定義
（定義~\ref{def:meas-weak-convergence}）をここで再掲すると，
\[
  \pi_k\Rightarrow\pi
  \quad\overset{\mathrm{def}}{\Longleftrightarrow}\quad
  \int F\,\d\pi_k\longrightarrow\int F\,\d\pi
  \qquad
  (\forall F\in C_b(\mathcal{X}\times\mathcal{X})).
\]

\begin{lemma}{$p$ 次輸送コストの下半連続性（Villani Lem 4.3）}{wp-cost-lsc}
  $1\leq p<\infty$ とする．
  $\Prob(\mathcal{X}\times\mathcal{X})$ の列 $(\pi_k)$ が
  $\pi\in\Prob(\mathcal{X}\times\mathcal{X})$ に弱収束，すなわち
  $\pi_k\Rightarrow\pi$ とする．このとき
  \[
    \int_{\mathcal{X}\times\mathcal{X}}d(x,y)^p\,\d\pi(x,y)
    \leq
    \liminf_{k\to\infty}
    \int_{\mathcal{X}\times\mathcal{X}}d(x,y)^p\,\d\pi_k(x,y).
  \]
\end{lemma}

\begin{proof}
  $d^p$ は連続だが有界とは限らないので，弱収束の定義をそのまま使うことは
  できない．そこで各 $m\in\N$ に対し
  \[
    c_m(x,y):=\min\{d(x,y)^p,m\}
  \]
  と切り捨てた関数を考え，$m\to\infty$ とする．
  $d$ は連続（注意~\ref{rem:wp-welldef}）だから $c_m$ も連続であり，
  $0\leq c_m\leq m$ より有界である．

  以下 $m$ を固定する．$c_m$ は有界連続だから，弱収束の定義
  （定義~\ref{def:meas-weak-convergence}）より
  \[
    \lim_{k\to\infty}\int c_m\,\d\pi_k=\int c_m\,\d\pi.
  \]
  一方 $c_m\leq d^p$ だから，各 $k$ で
  $\int c_m\,\d\pi_k\leq\int d^p\,\d\pi_k$ である．
  よって $\liminf$ の単調性（定義~\ref{def:meas-limsup}）を
  $a_k=\int c_m\,\d\pi_k$，$b_k=\int d^p\,\d\pi_k$ に適用すると，
  $(a_k)$ は収束するので $\lim_ka_k=\liminf_ka_k\leq\liminf_kb_k$，すなわち
  \[
    \int c_m\,\d\pi\leq\liminf_{k\to\infty}\int d^p\,\d\pi_k.
  \]

  最後に $m\to\infty$ とする．ここで用いる単調収束定理
  （定理~\ref{thm:meas-mct}）の主張は，測度空間
  $(\Omega,\mathcal{F},\lambda)$ 上の非負可測関数列 $(f_m)$ が

  \[
    f_m(\omega)\leq f_{m+1}(\omega),
    \qquad
    f_m(\omega)\longrightarrow f(\omega)
    \quad(\forall\omega\in\Omega)
  \]

  をみたすならば，

  \[
    \lim_{m\to\infty}\int_\Omega f_m\,\d\lambda
    =\int_\Omega f\,\d\lambda
  \]

  が成り立つ，というものである．積分値は $+\infty$ でもよい．
  今回は

  \[
    \Omega=\mathcal{X}\times\mathcal{X},
    \qquad
    \mathcal{F}=\Borel(\mathcal{X}\times\mathcal{X}),
    \qquad
    \lambda=\pi,
  \]
  \[
    f_m=c_m=\min\{d^p,m\},
    \qquad
    f=d^p
  \]

  として適用する．実際，$c_m$ は非負可測で，すべての $(x,y)$ に対して

  \[
    0\leq c_m(x,y)\leq c_{m+1}(x,y),
    \qquad
    c_m(x,y)\longrightarrow d(x,y)^p
  \]

  が成り立つ．したがって

  \[
    \lim_{m\to\infty}\int c_m\,\d\pi
    =\int d^p\,\d\pi.
  \]

  一方，先ほど得た不等式

  \[
    \int c_m\,\d\pi
    \leq
    \liminf_{k\to\infty}\int d^p\,\d\pi_k
  \]

  の右辺は $m$ によらない．よって両辺で $m\to\infty$ とし，
  上の単調収束定理の適用結果を用いると
  \[
    \int d^p\,\d\pi
    \leq\liminf_{k\to\infty}\int d^p\,\d\pi_k
  \]
  を得る．
\end{proof}

\begin{remark}{記法：測度の族}{wp-family-symbol}
  Villani は測度の族を $\mathcal{P},\mathcal{Q}$ と書くが，
  本資料では確率測度全体を $\Prob(\mathcal{X})$ と書いているため
  記号が衝突する．そこで以下では族を $A,B$ で表す．
\end{remark}

\begin{lemma}{輸送計画の tightness（Villani Lem 4.4）}{wp-plans-tight}
  $\mathcal{X},\mathcal{Y}$ を Polish 空間とし，
  $A\subset\Prob(\mathcal{X})$，$B\subset\Prob(\mathcal{Y})$ を
  tight（定義~\ref{def:meas-tight}）な確率測度の族とする．
  周辺分布がそれぞれ $A$，$B$ に属する coupling の全体を
  \[
    \Couplings(A,B)
    :=\bigcup_{\mu\in A,\,\nu\in B}\Couplings(\mu,\nu)
  \]
  と定めると，$\Couplings(A,B)$ は $\Prob(\mathcal{X}\times\mathcal{Y})$
  で tight である．
\end{lemma}

\begin{proof}
  $\varepsilon>0$ とする．$A,B$ の tightness より，compact 集合
  $K\subset\mathcal{X}$，$L\subset\mathcal{Y}$ で
  \[
    \eta(\mathcal{X}\setminus K)\leq\frac{\varepsilon}{2}
    \quad(\forall\eta\in A),
    \qquad
    \rho(\mathcal{Y}\setminus L)\leq\frac{\varepsilon}{2}
    \quad(\forall\rho\in B)
  \]
  をみたすものが存在する．$K,L$ は $\pi$ の選び方によらず，
  族 $A,B$ に対して一様に選ばれていることが重要である．

  ここで $\pi\in\Couplings(A,B)$ を任意に取る．
  $\Couplings(A,B)$ の定義より，

  \[
    \eta_\pi\in A,\qquad \rho_\pi\in B,\qquad
    \pi\in\Couplings(\eta_\pi,\rho_\pi)
  \]

  をみたす $\eta_\pi,\rho_\pi$ が存在する．上で選んだ $K,L$ に対して
  \[
    \eta_\pi(\mathcal{X}\setminus K)\leq\frac{\varepsilon}{2},
    \qquad
    \rho_\pi(\mathcal{Y}\setminus L)\leq\frac{\varepsilon}{2}.
  \]

  $\pi\in\Couplings(\eta_\pi,\rho_\pi)$ だから，
  \[
    \begin{aligned}
      \pi\bigl((\mathcal{X}\setminus K)\times\mathcal{Y}\bigr)
      &=\eta_\pi(\mathcal{X}\setminus K)
      \leq\frac{\varepsilon}{2},\\
      \pi\bigl(\mathcal{X}\times(\mathcal{Y}\setminus L)\bigr)
      &=\rho_\pi(\mathcal{Y}\setminus L)
      \leq\frac{\varepsilon}{2}.
    \end{aligned}
  \]

  集合の恒等式
  \[
    (\mathcal{X}\times\mathcal{Y})\setminus(K\times L)
    =\bigl((\mathcal{X}\setminus K)\times\mathcal{Y}\bigr)
      \cup
      \bigl(\mathcal{X}\times(\mathcal{Y}\setminus L)\bigr)
  \]
  と測度の劣加法性より，
  \[
    \begin{aligned}
      \pi\bigl((\mathcal{X}\times\mathcal{Y})
        \setminus(K\times L)\bigr)
      &\leq
      \pi\bigl((\mathcal{X}\setminus K)\times\mathcal{Y}\bigr)
      +\pi\bigl(\mathcal{X}\times(\mathcal{Y}\setminus L)\bigr)
      \\
      &\leq
      \frac{\varepsilon}{2}+\frac{\varepsilon}{2}\\
      &=\varepsilon.
    \end{aligned}
  \]
  $K,L$ は compact だから $K\times L$ も compact であり，
  上の評価は任意の $\pi\in\Couplings(A,B)$ に対して成り立つ．
  よって
  \[
    \pi\bigl((\mathcal{X}\times\mathcal{Y})\setminus(K\times L)\bigr)
    \leq\varepsilon
    \qquad(\forall\pi\in\Couplings(A,B)).
  \]
  定義~\ref{def:meas-tight}より $\Couplings(A,B)$ は tight である．
\end{proof}

\begin{theorem}{$\Wass_p$ の最適 coupling の存在（Villani Thm 4.1；pp.~43--45）}{wp-optimal-coupling-existence}
  $\mathcal{X}$ を Polish 空間，$1\leq p<\infty$ とする．
  任意の $\mu,\nu\in\Pp{p}(\mathcal{X})$ に対して，
  ある $\pi^*\in\Couplings(\mu,\nu)$ が存在し，
  \[
    \Wass_p(\mu,\nu)
    =
    \left(
      \int_{\mathcal{X}\times\mathcal{X}}
      d(x,y)^p\,\d\pi^*(x,y)
    \right)^{1/p}
  \]
  が成り立つ．
  すなわち，$\Wass_p$ の定義に現れる下限は最小値として達成される．
  この $\pi^*$ を $\Wass_p$ の\textbf{最適 coupling}という．
\end{theorem}

\begin{proof}
  各 $\pi\in\Couplings(\mu,\nu)$ に対して
  \[
    J(\pi)
    :=
    \int_{\mathcal{X}\times\mathcal{X}}d(x,y)^p\,\d\pi(x,y),
    \qquad
    I
    :=
    \inf_{\pi\in\Couplings(\mu,\nu)}J(\pi)
  \]
  とおく．$\Couplings(\mu,\nu)$ は空でない
  （補題~\ref{lem:wp-finite}の証明で $\mu\otimes\nu$ を構成した）．
  また $t\mapsto t^{1/p}$ は $[0,\infty]$ 上で単調増加だから，
  $\Wass_p$ の定義より
  \[
    \Wass_p(\mu,\nu)^p
    =
    \left(
      \inf_{\pi\in\Couplings(\mu,\nu)}J(\pi)^{1/p}
    \right)^p
    =
    \inf_{\pi\in\Couplings(\mu,\nu)}J(\pi)
    =I.
  \]
  補題~\ref{lem:wp-finite}より $\Wass_p(\mu,\nu)<\infty$ なので，
  \[
    0\leq I=\Wass_p(\mu,\nu)^p<\infty.
  \]

  \textbf{1. 最小化列を取る．}
  下限の定義より，各 $k\in\N$ に対して
  \[
    \pi_k\in\Couplings(\mu,\nu),
    \qquad
    I\leq J(\pi_k)<I+\frac1k
  \]
  となる $\pi_k$ を取ることができる．$I<\infty$ だから
  \[
    0\leq J(\pi_k)-I<\frac1k,
    \qquad
    J(\pi_k)\longrightarrow I.
  \]

  \textbf{2. 最小化列から弱収束部分列を取る．}
  一点族 $\{\mu\}$ の任意の列は定数列であり，$\mu$ に弱収束する．
  $\mathcal{X}$ は Polish 空間だから，Prokhorov の定理
  （定理~\ref{thm:meas-prokhorov}）より $\{\mu\}$ は tight である．
  同様に $\{\nu\}$ も tight である．
  したがって補題~\ref{lem:wp-plans-tight}より
  \[
    \Couplings(\mu,\nu)
    =\Couplings(\{\mu\},\{\nu\})
    \quad\text{は tight}.
  \]
  特に $(\pi_k)_{k\in\N}$ は tight である．
  $\mathcal{X}\times\mathcal{X}$ は Polish 空間なので，
  Prokhorov の定理より，ある部分列 $(\pi_{k_j})$ と
  $\pi^*\in\Prob(\mathcal{X}\times\mathcal{X})$ が存在して
  \[
    \pi_{k_j}\Rightarrow\pi^*
  \]
  が成り立つ．

  \textbf{3. $\pi^*$ も $\mu,\nu$ の coupling である．}
  任意の $\varphi\in C_b(\mathcal{X})$ に対し，
  \[
    F_\varphi(x,y):=\varphi(x)
  \]
  とおくと，$F_\varphi\in C_b(\mathcal{X}\times\mathcal{X})$ である．
  したがって $\pi_{k_j}\Rightarrow\pi^*$ と弱収束の定義より
  \[
    \int_{\mathcal{X}\times\mathcal{X}}
    \varphi(x)\,\d\pi_{k_j}(x,y)
    \longrightarrow
    \int_{\mathcal{X}\times\mathcal{X}}
    \varphi(x)\,\d\pi^*(x,y).
  \]
  一方，各 $\pi_{k_j}\in\Couplings(\mu,\nu)$ の第1周辺分布は
  $\mu$ だから，周辺積分公式（補題~\ref{lem:meas-marginal-integral}）より
  \[
    \int_{\mathcal{X}\times\mathcal{X}}
    \varphi(x)\,\d\pi_{k_j}(x,y)
    =
    \int_{\mathcal{X}}\varphi(x)\,\d\mu(x)
    \qquad(\forall j\in\N).
  \]
  よって
  \[
    \int_{\mathcal{X}\times\mathcal{X}}
    \varphi(x)\,\d\pi^*(x,y)
    =
    \int_{\mathcal{X}}\varphi(x)\,\d\mu(x).
  \]
  有界連続関数による測度の判定
  （定理~\ref{thm:meas-cb-determines}）より，
  $\pi^*$ の第1周辺分布は $\mu$ である．同様に，
  任意の $\psi\in C_b(\mathcal{X})$ に対して
  \[
    \int_{\mathcal{X}\times\mathcal{X}}
    \psi(y)\,\d\pi^*(x,y)
    =
    \int_{\mathcal{X}}\psi(y)\,\d\nu(y)
  \]
  が成り立つから，第2周辺分布は $\nu$ である．したがって
  \[
    \pi^*\in\Couplings(\mu,\nu).
  \]

  \textbf{4. $\pi^*$ が最小値を達成する．}
  $(x,y)\mapsto d(x,y)^p$ は非負連続関数なので，
  補題~\ref{lem:wp-cost-lsc}より
  \[
    J(\pi^*)
    \leq
    \liminf_{j\to\infty}J(\pi_{k_j}).
  \]
  一方，$\pi^*\in\Couplings(\mu,\nu)$ だから
  $I\leq J(\pi^*)$ である．また，
  $J(\pi_k)\to I$ より $J(\pi_{k_j})\to I$ である．したがって
  \[
    I
    \leq J(\pi^*)
    \leq\liminf_{j\to\infty}J(\pi_{k_j})
    =I.
  \]
  ゆえに $J(\pi^*)=I$ である．以上より
  \[
    \begin{aligned}
      \Wass_p(\mu,\nu)
      &=I^{1/p}\\
      &=J(\pi^*)^{1/p}\\
      &=
      \left(
        \int_{\mathcal{X}\times\mathcal{X}}
        d(x,y)^p\,\d\pi^*(x,y)
      \right)^{1/p}.
    \end{aligned}
  \]
  よって $\pi^*$ は $\Wass_p(\mu,\nu)$ の最適 coupling である．
\end{proof}

\section{Wasserstein 距離の距離性}
\label{sec:metric-property}

前節で最適 coupling の存在が分かったので，$\Wass_p$ が
$\Pp{p}(\mathcal{X})$ 上の距離の公理をみたすことを示す．
非負性と対称性は定義から直接従う．一致の公理の逆向き
$\Wass_p(\mu,\nu)=0\Rightarrow\mu=\nu$ では最適 coupling を用いる．
三角不等式では，$\mu$ から $\nu$ への輸送と $\nu$ から $\xi$ への輸送を
一つの三変数の確率測度にまとめるため，次の接着補題を用いる．

\begin{lemma}{接着補題（gluing lemma；Villani pp.~23--24）}{wp-gluing}
  $\mu,\nu,\xi\in\Prob(\mathcal{X})$，$\pi\in\Couplings(\mu,\nu)$，
  $\sigma\in\Couplings(\nu,\xi)$ とする
  （$\mathcal{X}^3$ には付録Aの直積距離を入れる）．
  このとき $\mathcal{X}^3$ 上の確率測度 $\gamma$ で，
  任意の $A,B,C\in\Borel(\mathcal{X})$ に対し
  \[
    \gamma(A\times B\times \mathcal{X})=\pi(A\times B),
    \qquad
    \gamma(\mathcal{X}\times B\times C)=\sigma(B\times C)
  \]
  をみたすものが存在する．
  すなわち，$\gamma$ の第 $1,2$ 成分の周辺分布が $\pi$ に，
  第 $2,3$ 成分の周辺分布が $\sigma$ に一致する．
\end{lemma}

\begin{proof}
  \textbf{1. 二つの coupling を共通の周辺分布 $\nu$ に沿って分解する．}

  測度分解定理（定理~\ref{thm:meas-disintegration}）を
  $\pi\in\Couplings(\mu,\nu)$ に適用すると，$\nu$ から $\mathcal{X}$ への
  Markov 核 $K_\pi$ が存在し，すべての
  $A,B\in\Borel(\mathcal{X})$ に対して
  \[
    \pi(A\times B)=\int_B K_\pi(y,A)\,\d\nu(y)
  \]
  が成り立つ．これを
  \[
    \pi(\d x,\d y)=K_\pi(y,\d x)\,\nu(\d y)
  \]
  と書く．

  $\sigma\in\Couplings(\nu,\xi)$ については第1，第2座標を入れ替えて
  同じ定理を適用する．すると $\nu$ から $\mathcal{X}$ への Markov 核
  $K_\sigma$ が存在して
  \[
    \sigma(\d y,\d z)=\nu(\d y)\,K_\sigma(y,\d z)
  \]
  と書ける．

  \textbf{2. 条件付き分布を使って三変数の測度を作る．}

  Markov 核による測度の貼り合わせ
  （定理~\ref{thm:meas-kernel-composition}）を
  $\nu,K_\pi,K_\sigma$ に適用すると，
  \[
    \gamma(\d x,\d y,\d z)
    :=K_\pi(y,\d x)\,\nu(\d y)\,K_\sigma(y,\d z)
  \]
  で表される $\gamma\in\Prob(\mathcal{X}^3)$ が存在する．具体的には，
  任意の $A,B,C\in\Borel(\mathcal{X})$ に対して
  \[
    \gamma(A\times B\times C)
    =\int_B K_\pi(y,A)K_\sigma(y,C)\,\d\nu(y)
  \]
  である．

  \textbf{3. 必要な二つの周辺分布を確かめる．}

  各 $y$ に対して $K_\pi(y,\mathcal{X})=K_\sigma(y,\mathcal{X})=1$
  だから，前式で $C=\mathcal{X}$，または $A=\mathcal{X}$ とおくと
  \[
    \begin{aligned}
      \gamma(A\times B\times\mathcal{X})
      &=\int_B K_\pi(y,A)K_\sigma(y,\mathcal{X})\,\d\nu(y)\\
      &=\int_B K_\pi(y,A)\,\d\nu(y)\\
      &=\pi(A\times B),\\[3pt]
      \gamma(\mathcal{X}\times B\times C)
      &=\int_B K_\pi(y,\mathcal{X})K_\sigma(y,C)\,\d\nu(y)\\
      &=\int_B K_\sigma(y,C)\,\d\nu(y)\\
      &=\sigma(B\times C).
    \end{aligned}
  \]
  したがって $\gamma$ は主張された二つの周辺分布をもつ．
\end{proof}

\begin{remark}{接着補題の確率的な意味}{wp-gluing-meaning}
  $\gamma$ は，まず $y\sim\nu$ を取り，$y$ を固定した条件のもとで
  $x\sim K_\pi(y,\cdot)$ と $z\sim K_\sigma(y,\cdot)$ を独立に取る，
  という三段階の操作で得られる同時分布である．
  したがって $(x,y)$ の同時分布は $\pi$，$(y,z)$ の同時分布は $\sigma$ のまま，
  二つの輸送を共通の中間点 $y$ を介して一つにまとめられる．
\end{remark}

\begin{proposition}{Wasserstein 距離は距離である（Villani Def 6.4）}{wp-metric}
  $1\leq p<\infty$ に対し，
  $\Wass_p$ は $\Pp{p}(\mathcal{X})$ 上の距離である．
\end{proposition}

\begin{proof}
  $\mu,\nu\in\Pp{p}(\mathcal{X})$ を任意に取る．
  距離の公理（定義~\ref{def:meas-metric-space}）を順に確かめる．

  \textbf{1. 有限性と非負性．}

  任意の $\pi\in\Couplings(\mu,\nu)$ に対して
  \[
    \left(\int d(x,y)^p\,\d\pi(x,y)\right)^{1/p}\geq0
  \]
  だから，その下限も $\Wass_p(\mu,\nu)\geq0$ である．
  また，補題~\ref{lem:wp-finite}より $\Wass_p(\mu,\nu)<\infty$ である．

  \textbf{2. 対称性．}

  $s(x,y):=(y,x)$ とおく．$\pi\in\Couplings(\mu,\nu)$ に対し
  像測度 $\pi':=\pi\circ s^{-1}$ を考える．任意の
  $A,B\in\Borel(\mathcal{X})$ に対して
  \[
    \begin{aligned}
      \pi'(A\times\mathcal{X})
      &=\pi(\mathcal{X}\times A)=\nu(A),\\
      \pi'(\mathcal{X}\times B)
      &=\pi(B\times\mathcal{X})=\mu(B).
    \end{aligned}
  \]
  よって $\pi'\in\Couplings(\nu,\mu)$ である．
  像測度の変数変換公式（補題~\ref{lem:meas-change-of-variables}）と
  距離の対称性 $d(y,x)=d(x,y)$ より
  \[
    \begin{aligned}
      \int d(x,y)^p\,\d\pi'(x,y)
      &=\int (d^p\circ s)(x,y)\,\d\pi(x,y)\\
      &=\int d(y,x)^p\,\d\pi(x,y)\\
      &=\int d(x,y)^p\,\d\pi(x,y).
    \end{aligned}
  \]
  $s\circ s=\Id$ だから $\pi\mapsto\pi\circ s^{-1}$ は
  $\Couplings(\mu,\nu)$ と $\Couplings(\nu,\mu)$ の間の全単射であり，
  対応する coupling の輸送コストは等しい．したがって両集合上で下限を取れば
  $\Wass_p(\mu,\nu)=\Wass_p(\nu,\mu)$ を得る．

  \textbf{3. 一致の公理：$\mu=\nu\Rightarrow\Wass_p(\mu,\nu)=0$．}

  対角写像 $\iota(x):=(x,x)$ と，その像測度
  $\pi_\Delta:=\mu\circ\iota^{-1}$ を考える．任意の
  $A,B\in\Borel(\mathcal{X})$ に対して
  \[
    \begin{aligned}
      \pi_\Delta(A\times\mathcal{X})
      &=\mu\bigl(\iota^{-1}(A\times\mathcal{X})\bigr)=\mu(A),\\
      \pi_\Delta(\mathcal{X}\times B)
      &=\mu\bigl(\iota^{-1}(\mathcal{X}\times B)\bigr)=\mu(B).
    \end{aligned}
  \]
  よって $\pi_\Delta\in\Couplings(\mu,\mu)$ である．さらに，
  像測度の積分公式と $d(x,x)=0$ より
  \[
    \int d(x,y)^p\,\d\pi_\Delta(x,y)
    =\int_\mathcal{X}d(x,x)^p\,\d\mu(x)
    =0.
  \]
  $\Wass_p$ は非負であり，コスト $0$ の coupling が存在するから
  $\Wass_p(\mu,\mu)=0$ である．

  \textbf{4. 一致の公理：$\Wass_p(\mu,\nu)=0\Rightarrow\mu=\nu$．}

  $\Wass_p(\mu,\nu)=0$ と仮定する．最適 coupling の存在定理
  （定理~\ref{thm:wp-optimal-coupling-existence}）より，ある
  $\pi^*\in\Couplings(\mu,\nu)$ が存在して
  \[
    0=\Wass_p(\mu,\nu)^p
    =\int d(x,y)^p\,\d\pi^*(x,y)
  \]
  が成り立つ．被積分関数 $d^p$ は非負だから，積分が $0$ であることから
  $d(x,y)^p=0$，したがって $d(x,y)=0$ が $\pi^*$-a.e. で成り立つ．
  距離の一致の公理より，これは $x=y$ が $\pi^*$-a.e. であることを意味する．

  対角集合を $\Delta:=\{(x,x)\mid x\in\mathcal{X}\}$ とおく．
  $\Delta=\{(x,y)\mid d(x,y)=0\}$ は連続関数 $d$ による閉集合
  $\{0\}$ の逆像なので閉集合，特に Borel 集合であり，
  $\pi^*(\Delta)=1$ である．したがって任意の $A\in\Borel(\mathcal{X})$ に対し
  \[
    \mu(A)
    =\pi^*(A\times \mathcal{X})
    =\pi^*((A\times\mathcal{X})\cap\Delta)
    =\pi^*((\mathcal{X}\times A)\cap\Delta)
    =\pi^*(\mathcal{X}\times A)
    =\nu(A).
  \]
  よって $\mu=\nu$ である．

  \textbf{5. 三角不等式．}

  $\xi\in\Pp{p}(\mathcal{X})$ を任意に取る．最適 coupling の存在定理より，
  最適 coupling
  $\pi\in\Couplings(\mu,\nu)$，$\sigma\in\Couplings(\nu,\xi)$ を
  それぞれ取ることができる．したがって
  \[
    \begin{aligned}
      \Wass_p(\mu,\nu)
      &=\left(\int d(x,y)^p\,\d\pi(x,y)\right)^{1/p},\\
      \Wass_p(\nu,\xi)
      &=\left(\int d(y,z)^p\,\d\sigma(y,z)\right)^{1/p}.
    \end{aligned}
  \]

  接着補題（補題~\ref{lem:wp-gluing}）により，
  $\gamma\in\Prob(\mathcal{X}^3)$ で
  \[
    \gamma(A\times B\times\mathcal{X})=\pi(A\times B),
    \qquad
    \gamma(\mathcal{X}\times B\times C)=\sigma(B\times C)
  \]
  （$A,B,C\in\Borel(\mathcal{X})$）となるものがある．

  $\gamma$ の第1，第3成分の周辺分布を
  $\tau\in\Prob(\mathcal{X}\times\mathcal{X})$ とする．
  任意の $A,C\in\Borel(\mathcal{X})$ に対して
  \[
    \begin{aligned}
      \tau(A\times\mathcal{X})
      &=\gamma(A\times\mathcal{X}\times\mathcal{X})\\
      &=\pi(A\times\mathcal{X})=\mu(A),\\
      \tau(\mathcal{X}\times C)
      &=\gamma(\mathcal{X}\times\mathcal{X}\times C)\\
      &=\sigma(\mathcal{X}\times C)=\xi(C).
    \end{aligned}
  \]
  したがって $\tau\in\Couplings(\mu,\xi)$ である．

  距離 $d$ の三角不等式より，すべての $(x,y,z)\in\mathcal{X}^3$ で
  \[
    0\leq d(x,z)\leq d(x,y)+d(y,z).
  \]
  両辺の $L^p(\gamma)$ ノルムを取り，$L^p$ ノルムの単調性
  （主張~\ref{clm:meas-lp-monotone}）と Minkowski の不等式
  （定理~\ref{thm:meas-minkowski}）を順に用いると
  \[
    \norm{d(x,z)}_{L^p(\gamma)}
    \leq
    \norm{d(x,y)}_{L^p(\gamma)}
    +
    \norm{d(y,z)}_{L^p(\gamma)}.
  \]
  ここで $\gamma$ の第1，第2成分の周辺分布は $\pi$，
  第2，第3成分の周辺分布は $\sigma$ だから，周辺積分公式
  （補題~\ref{lem:meas-marginal-integral}）より
  \[
    \begin{aligned}
      \norm{d(x,y)}_{L^p(\gamma)}^p
      &=\int_{\mathcal{X}^3}d(x,y)^p\,\d\gamma(x,y,z)\\
      &=\int_{\mathcal{X}\times\mathcal{X}}d(x,y)^p\,\d\pi(x,y),\\
      \norm{d(y,z)}_{L^p(\gamma)}^p
      &=\int_{\mathcal{X}^3}d(y,z)^p\,\d\gamma(x,y,z)\\
      &=\int_{\mathcal{X}\times\mathcal{X}}d(y,z)^p\,\d\sigma(y,z).
    \end{aligned}
  \]
  よって $\pi,\sigma$ の最適性から
  \[
    \norm{d(x,y)}_{L^p(\gamma)}=\Wass_p(\mu,\nu),
    \qquad
    \norm{d(y,z)}_{L^p(\gamma)}=\Wass_p(\nu,\xi).
  \]

  一方，$\tau$ は $\mu,\xi$ の coupling だから，$\Wass_p$ の定義より
  \[
    \Wass_p(\mu,\xi)
    \leq
    \left(
      \int_{\mathcal{X}\times\mathcal{X}}d(x,z)^p\,\d\tau(x,z)
    \right)^{1/p}.
  \]
  $\tau$ は $\gamma$ の第1，第3成分の周辺分布なので，
  再び周辺積分公式を用いると，右辺は
  $\norm{d(x,z)}_{L^p(\gamma)}$ に等しい．以上をまとめて
  \[
    \begin{aligned}
      \Wass_p(\mu,\xi)
      &\leq\norm{d(x,z)}_{L^p(\gamma)}\\
      &\leq
      \norm{d(x,y)}_{L^p(\gamma)}
      +\norm{d(y,z)}_{L^p(\gamma)}\\
      &=\Wass_p(\mu,\nu)+\Wass_p(\nu,\xi)
    \end{aligned}
  \]
  を得る．これは $\Wass_p$ の三角不等式である．

  以上で距離の公理がすべて示された．
\end{proof}

\begin{definition}{Wasserstein 空間（Villani Def 6.4）}{wp-wasserstein-space}
  $1\leq p<\infty$ に対し，距離 $\Wass_p$ を備えた
  $\Pp{p}(\mathcal{X})$ を $p$ 次の \textbf{Wasserstein 空間}といい，
  \[
    (\Pp{p}(\mathcal{X}),\Wass_p)
  \]
  と書く．
\end{definition}

\begin{remark}{文献との対応}{wp-metric-references}
  存在定理は Villani (2009) Theorem 4.1 をコスト $d^p$ に特殊化したものであり，
  二つの補題は同 Lemmas 4.3, 4.4 に対応する．
  距離性の証明は同 Chapter 6 に従う．接着補題は同 Chapter 1 の
  Gluing Lemma（桑江ほか (2015) 補題 1.33）である．
  なお Villani (2009) Definition 6.1, 6.4 は
  まず $\Prob(\mathcal{X})$ 上に $\Wass_p$ を定義し，
  その後 $\Pp{p}(\mathcal{X})$ に制限して有限値の距離とするが，
  本資料でははじめから $\Pp{p}(\mathcal{X})$ 上で定義した．
\end{remark}
